% SEGMENT OLP-0639-S01
% Part: history
% Chapter: set-theory
% Section: hilbert-curve

\documentclass[../../../include/open-logic-section]{subfiles}

\begin{document}

\olfileid{his}{set}{hilbertcurve}

\olsection{இணைப்பு: ஹில்பர்ட்டின் வெளிநிரப்பு வளைவுகள்}

\olref[pathology]{sec} அத்தியாயத்தில், ஒரு கோட்டிலும்
ஒரு சதுரத்திலும் ஒரே எண்ணிக்கையிலான புள்ளிகள் உள்ளன என்ற
காண்டோரின் நிறுவல்
(\olref[his][set][cantorplane]{thm:cantorplane}),
வெளியை நிரப்பும் \emph{வளைவு} இருக்குமா என்று பியானோவைக்
கேட்கத் தூண்டியதைக் குறிப்பிட்டோம். 1890-இல் அதற்கு
உடன்பாடான பதிலை அவர் பெற்றார். இப்பகுதியில், சிறிய மாற்றத்துடன்,
ஹில்பர்ட்டின் வெளிநிரப்பு வளைவை எவ்வாறு கட்டலாம் என்பதை
(முழு முறைசார்ந்த நிறுவலாக அல்லாமல்) விளக்குகிறோம்.\footnote{இன்னும்
துல்லியமான விளக்கத்திற்கு \cite{Rose2010}-ஐப் பார்க்கவும்.
கீழுள்ள வளைவுகளின் சிவப்புப் பகுதிகளைச் சேர்ப்பதே இங்கு
செய்யப்படும் சிறிய மாற்றம். இதனால் வளைவு தொடர்ச்சியானது
என்பதைச் சரிபார்ப்பது சற்றே எளிதாகிறது.}

$h_1$, $h_2$, $h_3$, \dots\@ என்ற சார்புகளின் தொடரின்
எல்லையாக $h$ என்ற சார்பை வரையறுக்க வேண்டும். முதலில்
கட்டுமானத்தை விளக்குகிறோம். பின்னர் அது வெளியை நிரப்புகிறது
என்று காட்டுகிறோம். இறுதியாக அது ஒரு வளைவு என்றும்
காட்டுகிறோம்.

% SEGMENT OLP-0639-S02
$h$-இன் வீச்சாக அலகுச் சதுரமான $\unitsquare$-ஐ
எடுத்துக்கொள்வோம். $h$-க்கான முதல் தோராயம், அதாவது
$h_1$, இதோ:
\begin{center}
	\begin{tikzpicture}
	\draw[gray] (0pt, 0pt) rectangle (80pt, 80pt);
	\draw[step = 40pt, gray, very thin] (0pt, 0pt) grid (80pt, 80pt);
	\draw[oldiagcolorC, thick] (20pt, 0)--(20pt, 20pt);
	\draw[oldiagcolorC, thick] (60pt, 0)--(60pt, 20pt);
	\begin{scope}[xshift=20pt, yshift=20pt]
	\draw [black, thick, l-system={Hilbert curve, step=40pt, angle=90, axiom=L, order=1}]  lindenmayer system;
	\end{scope}
	\end{tikzpicture}
\end{center}
நிலைகளைப் பின்தொடர உதவ, சதுரத்தின்மேல் $2 \times 2$
கட்டத்தை அமைத்திருக்கிறோம். வளைவு கீழ் இடது கால்
பகுதியில் தொடங்கி, மேல் இடது பகுதிக்கும், பின்னர் மேல்
வலது பகுதிக்கும், இறுதியாகக் கீழ் வலது பகுதிக்கும்
செல்வதாகக் கருதலாம். கட்டுமானத்தின் இரண்டாம் நிலை,
அதாவது $h_2$, இதோ:
\begin{center}
	\begin{tikzpicture}
	\draw[gray] (0pt, 0pt) rectangle (80pt, 80pt);
	\draw[step = 20pt, gray, very thin] (0pt, 0pt) grid (80pt, 80pt);
	\draw[oldiagcolorC, thick] (0pt, 10pt)--(10pt, 10pt);
	\draw[oldiagcolorC, thick] (70pt, 10pt)--(80pt, 10pt);
	\draw[thick, oldiagcolorE] (10pt, 30pt)--(10pt, 50pt);
	\draw[thick, oldiagcolorE] (30pt, 50pt)--(50pt, 50pt);
	\draw[thick, oldiagcolorE] (70pt, 50pt)--(70pt, 30pt); \begin{scope}[xshift=10pt, yshift=30pt]
	\draw [thick, rotate=270,black, l-system={Hilbert curve, step=20pt, angle=90, axiom=L, order=1}]  lindenmayer system;
	\end{scope}
	\begin{scope}[xshift=10pt, yshift=50pt]
	\draw [thick, black, l-system={Hilbert curve, step=20pt, angle=90, axiom=L, order=1}]  lindenmayer system;
	\end{scope}
	\begin{scope}[xshift=50pt, yshift=50pt]
	\draw [thick, black, l-system={Hilbert curve, step=20pt, angle=90, axiom=L, order=1}]  lindenmayer system;
	\end{scope}
	\begin{scope}[xshift=70pt, yshift=10pt]
	\draw [thick, rotate=90,black, l-system={Hilbert curve, step=20pt, angle=90, axiom=L, order=1}]  lindenmayer system;
	\end{scope}
	\end{tikzpicture}
\end{center}

% SEGMENT OLP-0639-S03
வெவ்வேறு நிறங்கள், $h_2$ எப்படி அமைக்கப்பட்டது
என்பதை விளக்க உதவும். முதலில் $h_1$-இன் !!{colorC}
அல்லாத பகுதியின் அளவு குறைக்கப்பட்ட நகல்களைச்
சதுரத்தின் கீழ் இடது, மேல் இடது, மேல் வலது, கீழ்
வலது பகுதிகளில் வைக்கிறோம் (கருப்பு நிறத்தில்).
பிறகு இந்த நான்கு வடிவங்களையும் (!!{colorE} கோடுகளால்)
இணைக்கிறோம். இறுதியாக வடிவத்தைச் சதுரத்தின்
எல்லையுடன் (!!{colorC} கோடுகளால்) இணைக்கிறோம்.
% SEGMENT OLP-0639-S04
இப்போது $h_3$-ஐப் பார்ப்போம். $h_2$ ஆனது
$h_1$-இன் சிவப்பு அல்லாத பகுதியின் அளவு குறைக்கப்பட்ட
நான்கு நகல்களை இணைத்து அமைக்கப்பட்டதுபோல,
$h_3$ ஆனது $h_2$-இன் சிவப்பு அல்லாத பகுதியின்
அளவு குறைக்கப்பட்ட நான்கு நகல்களால் (கருப்பு நிறத்தில்)
அமைக்கப்படுகிறது. அவை பிறகு (!!{colorE} கோடுகளால்)
ஒன்றோடொன்று இணைக்கப்பட்டு, இறுதியாக
(!!{colorC} கோடுகளால்) சதுரத்தின் எல்லையுடன்
இணைக்கப்படுகின்றன.
\begin{center}
	\begin{tikzpicture}
	\draw[gray] (0pt, 0pt) rectangle (80pt, 80pt);
	\draw[step = 10pt, gray, very thin] (0pt, 0pt) grid (80pt, 80pt);
	\draw[thick, oldiagcolorC](5pt, 0pt)--(5pt, 5pt);
	\draw[thick, oldiagcolorC] (75pt, 0pt)--(75pt, 5pt);
	\draw[thick, oldiagcolorE] (75pt, 45pt)--(75pt, 35pt);
	\draw[thick, oldiagcolorE] (5pt, 35pt)--(5pt, 45pt);
	\draw[thick, oldiagcolorE] (35pt, 45pt)--(45pt, 45pt);
	\draw[thick, oldiagcolorE] (75pt, 45pt)--(75pt, 35pt); \begin{scope}[xshift=5pt, yshift=35pt]
	\draw [rotate=270, thick, black, l-system={Hilbert curve, step=10pt, angle=90, axiom=L, order=2}]  lindenmayer system;
	\end{scope}
	\begin{scope}[xshift=5pt, yshift=45pt]
	\draw [thick, black, l-system={Hilbert curve, step=10pt, angle=90, axiom=L, order=2}]  lindenmayer system;
	\end{scope}
	\begin{scope}[xshift=45pt, yshift=45pt]
	\draw [thick, black, l-system={Hilbert curve, step=10pt, angle=90, axiom=L, order=2}]  lindenmayer system;
	\end{scope}
	\begin{scope}[xshift=75pt, yshift=5pt]
	\draw [thick, rotate=90,black, l-system={Hilbert curve, step=10pt, angle=90, axiom=L, order=2}]  lindenmayer system;
	\end{scope}
	\end{tikzpicture}
\end{center}

% SEGMENT OLP-0639-S05
இப்போது $h_n$-இலிருந்து $h_{n+1}$-ஐ
வரையறுக்கும் பொதுவான முறையைக் காணலாம்.
இறுதியாக, அடுத்தடுத்த சார்புகளான $h_1$, $h_2$,
\dots\@ ஆகியவற்றின் புள்ளிவாரி எல்லையைக் கொண்டு
வளைவாகிய $h$-ஐ \emph{அதையே} வரையறுக்கிறோம்.
அதாவது, ஒவ்வொரு $x \in \unitline$-க்கும்:
% SOURCE CORRECTION TA-HIS-004: the curve parameter lies in the unit line, not the unit square.
\begin{align*}
	h(x) &= \lim_{n \rightarrow \infty} h_n(x)
\end{align*}

% SEGMENT OLP-0639-S06
இப்போது இந்த வளைவு வெளியை நிரப்புகிறது என்று
காட்டுவோம். $h_n$ வளைவை வரையும்போது,
$\unitsquare$ மீது $2^n \times 2^n$ கட்டத்தை
அமைக்கிறோம். பித்தகோரசின் தேற்றப்படி, ஒவ்வொரு
கட்டச் சிறுபகுதியின் மூலைவிட்ட நீளம்:
\[
\sqrt{\left(\nicefrac{1}{2^{n}}\right)^2+\left(\nicefrac{1}{2^{n}}\right)^2} = 2^{(\frac{1}{2}-n)}
\]
மேலும், $h_n$ ஒவ்வொரு கட்டச் சிறுபகுதியையும்
கடந்து செல்கிறது. எனவே $\unitsquare$-இல்
உள்ள ஒவ்வொரு புள்ளியும் $h_n$-இல் உள்ள ஏதோ
ஒரு புள்ளியிலிருந்து \emph{அதிகபட்சம்}
$2^{(\frac{1}{2}-n)}$ தூரத்தில்தான் உள்ளது.
$h$ என்பது $h_1$, $h_2$, $h_3$, \dots\@
என்ற சார்புகளின் எல்லையாக வரையறுக்கப்பட்டது.
எனவே எந்தப் புள்ளிக்கும் $h$-க்கும் இடையிலான
அதிகபட்சத் தூரம்:
\[
\lim_{n \rightarrow \infty} 2^{(\frac{1}{2}-n)} = 0.
\]
அதாவது, $\unitsquare$-இல் உள்ள ஒவ்வொரு
புள்ளியும் $h$-இலிருந்து $0$ தூரத்தில் உள்ளது.
வேறு விதமாகச் சொன்னால், $\unitsquare$-இன்
ஒவ்வொரு புள்ளியும் வளைவின் \emph{மேலேயே}
உள்ளது. ஆகவே $h$ வெளியை நிரப்புகிறது!{}

% SEGMENT OLP-0639-S07
$h$ உண்மையில் ஒரு \emph{வளைவு} என்பதைக்
காட்டுவது மட்டுமே மீதமுள்ளது. இதற்காக அந்தக்
கருத்தை வரையறுக்க வேண்டும். நவீன வரையறை,
1887-இல் ஜோர்டான் அளித்த வரையறையை அடிப்படையாகக்
கொண்டது (அதாவது, முதல் வெளிநிரப்பு வளைவு
வழங்கப்படுவதற்கு சில ஆண்டுகளுக்கு முன்பு):

\begin{defn}
ஒரு வளைவு என்பது $\unitline$-இலிருந்து
$\Real^2$-க்குச் செல்லும் தொடர்ச்சியான படமாக்கல்.
\end{defn}

இது உள்ளுணர்வுக்கு மிகவும் பொருந்துகிறது:
ஒரு வளைவு என்பது, வழக்கமான ஒரு கோட்டிலிருந்து
$\Real^2$ தளத்திற்குச் செல்லும் “சீரான”
படமாக்கல் என்று கருதலாம். நமது $h$ உண்மையில்
$\unitline$-இலிருந்து $\Real^2$-க்குச்
செல்லும் படமாக்கல். எனவே $h$ தொடர்ச்சியானது
என்று காட்டினால் போதும். தொடர்ச்சியை
\olref[limits]{sec} பகுதியில் $\epsilon$/$\delta$
குறியீட்டால் வரையறுத்தோம். அன்றாடச் சொற்களில்,
பின்வருவதை நிறுவ விரும்புகிறோம்:
\emph{$\unitsquare$-இல் ஒரு புள்ளி $p$-யையும்,
விரும்பிய துல்லிய அளவு $\epsilon$-ஐயும்
குறிப்பிட்டால், $\unitline$-இல் ஒரு திறந்த
இடைவெளியைக் காணலாம்; அந்த இடைவெளியில் உள்ள
எந்த $x$-ஐ எடுத்தாலும் $h(x)$ ஆனது $p$-இலிருந்து
$\epsilon$ தூரத்திற்குள் இருக்கும்.}

% SEGMENT OLP-0639-S08
ஆகவே $p$-யையும் $\epsilon$-ஐயும்
குறிப்பிட்டிருக்கிறீர்கள் என்று கொள்வோம்.
இது சாராம்சத்தில், $\unitsquare$-இல் $p$-ஐ
மையமாகவும் $\epsilon$-ஐ ஆரமாகவும் கொண்ட
ஒரு வட்டத்தை வரைவதாகும். (வட்டத்தின் ஒரு பகுதி
$\unitsquare$-இன் எல்லையைத் தாண்டிச்
செல்லலாம்; அது பொருட்டல்ல.) $h_n$ சார்பை
விளக்கும்போது $\unitsquare$ மீது
$2^n \times 2^n$ கட்டத்தை வரைந்ததை
நினைவுகூருங்கள். $\epsilon$ எவ்வளவு
சிறியதாக இருந்தாலும், ஏதோ ஓர் $n$-க்கு
$\unitsquare$ மீது அமைந்த $2^n \times 2^n$
கட்டத்தின் ஒரு சிறுபகுதி முழுவதும்,
$p$ மையமும் $\epsilon$ ஆரமும் கொண்ட
வட்டத்திற்குள் அமையும் என்பது தெளிவு.

% SEGMENT OLP-0639-S09
அந்த $n$-ஐ எடுத்துக்கொண்டு, $h_n$ முழுவதுமாக
அந்தக் கட்டச் சிறுபகுதிக்குள் படமாக்கும்
$\unitline$-இன் மிகப் பெரிய திறந்த
பகுதியை $I$ என்று கொள்வோம். (அத்தகைய
$(a,b)$ இருப்பது தெளிவு; ஏனெனில் $h_n$,
$2^n\times 2^n$ கட்டத்தின் ஒவ்வொரு
சிறுபகுதியையும் கடந்து செல்கிறது என்று
ஏற்கெனவே குறிப்பிட்டோம்.) இப்போது
$x \in I$ எனும்போதெல்லாம் $h(x)$-உம்
அதே கட்டச் சிறுபகுதியில் இருப்பதைக்
காட்டினால் போதும். இதைச் செய்ய,
$m > n$ என்ற எந்த மதிப்புக்கும்
$h_m(x)$ அதே கட்டச் சிறுபகுதியில்
இருப்பதைக் காட்டினால் போதும். \emph{இதுவும்}
தெளிவு. $m > n$ எனும்போது $h_m$-இல்
என்ன நடக்கிறது என்று பார்த்தால்,
அலகு இடைவெளியின் துல்லியமாக “அதே பகுதி”
அதே கட்டச் சிறுபகுதிக்குள் படமாக்கப்படுகிறது;
அந்தப் பகுதியில் வளைவு மேலும் மேலும்
நீண்டும் வளைந்தும் செல்கிறது.

% SOURCE CORRECTION TA-HIS-005: the source's final neighbourhood argument is insufficient for continuity.
பதிப்பாசிரியர் குறிப்பு: மேலுள்ள கட்டுமான விளக்கம்
முறைசாரா ஒன்று. கடைசி வாதம், கொடுக்கப்பட்ட
வெளியீட்டுப் புள்ளியின் அருகில் படமாகும்
ஓர் உள்ளீட்டு இடைவெளியைத் தேடுகிறது; அது
குறிப்பிட்ட எந்த உள்ளீட்டுப் புள்ளியிலும்
எல்லைச் சார்பு தொடர்ச்சியானது என்பதை நிறுவவில்லை.
மேலும், புள்ளிவாரி எல்லையை மட்டுமே
இங்கு கூறியிருப்பதால், எல்லை வளைவு
முழுச் சதுரத்தையும் நிரப்புகிறது என்ற
முடிவிற்குத் தேவையான அணுகல் பண்பும்
இங்கே நிறுவப்படவில்லை. எனவே இந்தப்
பத்திகளை முழு நிறுவலாகக் கொள்ளக் கூடாது.

\end{document}
